\documentclass{article}
\usepackage{amsmath,amssymb,amsthm}

\theoremstyle{definition}
\newtheorem{definition}{Definition}

\theoremstyle{plain}
\newtheorem{lemma}{Lemma}
\newtheorem{theorem}{Theorem}

\begin{document}

\begin{definition}[Planning task]
A STRIPS planning task over a finite set of objects $O$ is a tuple $P = \langle S,A,I,G\rangle$. Its atoms are the ground expressions $p(o_1,\dots,o_k)$ for a predicate $p$ of arity $k$ and $o_i \in O$. A state $s \in S$ is a set of atoms, $I \in S$ is the initial state, $G$ is a set of atoms, and every action $a \in A$ is a triple $\langle \mathrm{pre}(a), \mathrm{add}(a), \mathrm{del}(a) \rangle$ of sets of atoms.
\end{definition}

\begin{definition}[Applicability and progression]\label{def:progression}
An action $a \in A$ is \emph{applicable} in a state $s \in S$ if $\mathrm{pre}(a) \subseteq s$, and the state it leads to is then
\[
  s[a] = \bigl(s \setminus \mathrm{del}(a)\bigr) \cup \mathrm{add}(a).
\]
A plan is a finite sequence $\pi = a_1 \dots a_n$ of actions. The empty plan $\varepsilon$ is applicable in every state, with $s[\varepsilon] = s$, and $\pi = a_1 \dots a_n$ with $n \geq 1$ is \emph{applicable} in $s$ if $a_1$ is applicable in $s$ and $a_2 \dots a_n$ is applicable in $s[a_1]$. The plan then leads to $s[\pi] = s[a_1][a_2 \dots a_n]$. A plan $\pi$ \emph{solves} $P$ if it is applicable in $I$ and $G \subseteq I[\pi]$.
\end{definition}

\begin{definition}[Object collapse]
Fix a set $C \subseteq O$ of objects to collapse, with $|C| > 1$, and a fresh symbol $c \notin O$. The collapse is the map $\sigma \colon O \to (O \setminus C) \cup \{c\}$ given by
\[
  \sigma(o) =
  \begin{cases}
    c & \text{if } o \in C,\\
    o & \text{otherwise.}
  \end{cases}
\]
\end{definition}

\begin{definition}[Abstraction map]
The abstraction map $(\cdot)^*$ lifts $\sigma$ to atoms, states, actions and plans by setting
\[
  p(o_1,\dots,o_k)^* = p(\sigma(o_1),\dots,\sigma(o_k)),
  \qquad
  s^* = \{\, f^* : f \in s \,\},
\]
\[
  a^* = \langle \mathrm{pre}(a)^*, \mathrm{add}(a)^*, \mathrm{del}(a)^* \rangle,
  \qquad
  (a_1 \dots a_n)^* = a_1^* \dots a_n^*.
\]
\end{definition}

\begin{definition}[Abstract task]
The abstraction of $P$ under $\sigma$ is $P^* = \langle S^*,A^*,I^*,G^*\rangle$, where $S^*$ collects the states over $(O \setminus C) \cup \{c\}$, $A^* = \{\, a^* : a \in A \,\}$, and $(\cdot)^*$ carries $I$ and $G$ to $I^*$ and $G^*$.
\end{definition}

\begin{definition}[Relaxed abstract task]\label{def:relaxed}
Say that an atom $p(o_1,\dots,o_k)$ \emph{mentions} an object $o$ if $o \in \{o_1,\dots,o_k\}$. For $a \in A$ put
\[
  \mathrm{del}^r(a) = \{\, d^* : d \in \mathrm{del}(a)
    \text{ and } d \text{ mentions no object of } C \,\},
\]
\[
  a^r = \langle\, \mathrm{pre}(a)^*,\ \mathrm{add}(a)^*,\ \mathrm{del}^r(a) \,\rangle,
\]
so that $a^r$ forgets every delete whose atom mentions a collapsed object. The \emph{relaxed abstract task} is $P^r = \langle S^*, A^r, I^*, G^* \rangle$ with $A^r = \{\, a^r : a \in A \,\}$. It differs from $P^*$ only in the delete lists of its actions, and coincides with $P^*$ when no delete of $P$ mentions an object of $C$.
\end{definition}

Since $\sigma$ maps every object of $C$ onto the single symbol $c$ it is not injective, so $(\cdot)^*$ may identify distinct atoms, states and actions of $P$.

\begin{lemma}\label{lem:action}
If $a \in A$ is applicable in $s \in S$, then $a^r$ is applicable in $s^*$.
\end{lemma}

\begin{proof}
By Definition~\ref{def:relaxed} the relaxation leaves preconditions untouched, so $\mathrm{pre}(a^r) = \mathrm{pre}(a)^*$. From $\mathrm{pre}(a) \subseteq s$ we get $\mathrm{pre}(a)^* \subseteq s^*$, hence $\mathrm{pre}(a^r) \subseteq s^*$.
\end{proof}

\begin{lemma}[Progression]\label{lem:progress}
If $a \in A$ is applicable in $s \in S$, then $(s[a])^* \subseteq s^*[a^r]$.
\end{lemma}

\begin{proof}
Progression gives $s[a] = (s \setminus \mathrm{del}(a)) \cup \mathrm{add}(a)$, so
\[
  (s[a])^* = (s \setminus \mathrm{del}(a))^* \cup \mathrm{add}(a)^*,
  \qquad
  s^*[a^r] = (s^* \setminus \mathrm{del}^r(a)) \cup \mathrm{add}(a)^*,
\]
and it suffices to show that $(s \setminus \mathrm{del}(a))^* \subseteq s^* \setminus \mathrm{del}^r(a)$. Let $f \in (s \setminus \mathrm{del}(a))^*$. Then $f = g^*$ for some $g \in s$ such that $g \notin \mathrm{del}(a)$, so $f \in s^*$. Suppose $f \in \mathrm{del}^r(a)$. By Definition~\ref{def:relaxed} we then have $f = d^*$ for some $d \in \mathrm{del}(a)$ mentioning no object of $C$. Since $\sigma$ fixes every object outside $C$, it fixes $d$, so $f = d$. Now $d$ is an atom over $O$ and $c \notin O$, so $f$ does not mention $c$. If $g$ mentioned an object of $C$ then $g^*$ would mention $c$, and $g^* = f$ does not. Hence $g$ mentions no object of $C$ either, so $\sigma$ fixes $g$ and $g = g^* = f = d$. This contradicts $g \notin \mathrm{del}(a)$, so $f \notin \mathrm{del}^r(a)$.
\end{proof}

\begin{lemma}[Monotonicity]\label{lem:mono}
Let $s \subseteq t$ be states and let $\pi$ be applicable in $s$. Then $\pi$ is applicable in $t$ and $s[\pi] \subseteq t[\pi]$.
\end{lemma}

\begin{proof}
By induction on the length of $\pi$. The empty plan is applicable in every state and $s[\varepsilon] = s \subseteq t = t[\varepsilon]$. For $\pi = a\pi'$, applicability of $a$ in $s$ gives $\mathrm{pre}(a) \subseteq s \subseteq t$, so $a$ is applicable in $t$, and
\[
  s[a] = (s \setminus \mathrm{del}(a)) \cup \mathrm{add}(a)
           \subseteq (t \setminus \mathrm{del}(a)) \cup \mathrm{add}(a)
            = t[a].
\]
The claim follows by applying the induction hypothesis to $\pi'$ at $s[a] \subseteq t[a]$.
\end{proof}

\begin{lemma}\label{lem:plan}
If $\pi$ is applicable in $s \in S$, then $\pi^r$ is applicable in $s^*$ and $(s[\pi])^* \subseteq s^*[\pi^r]$.
\end{lemma}

\begin{proof}
By induction on the length of $\pi$. For the empty plan both sides are $s^*$. Let $\pi = a\pi'$ with $a$ applicable in $s$. By Lemma~\ref{lem:action} $a^r$ is applicable in $s^*$, and by Lemma~\ref{lem:progress} $(s[a])^* \subseteq s^*[a^r]$. The induction hypothesis applied to $\pi'$ at $s[a]$ gives that $\pi'^r$ is applicable in $(s[a])^*$ with $(s[a][\pi'])^* \subseteq (s[a])^*[\pi'^r]$. Since $(s[a])^* \subseteq s^*[a^r]$, Lemma~\ref{lem:mono} lets us move both to $s^*[a^r]$, where $\pi'^r$ is applicable and
\[
  (s[\pi])^* = (s[a][\pi'])^* \subseteq (s[a])^*[\pi'^r]
             \subseteq s^*[a^r][\pi'^r] = s^*[\pi^r]. \qedhere
\]
\end{proof}

\begin{theorem}[Over-approximation]\label{thm:goal}
If $G \subseteq I[\pi]$, then $G^* \subseteq I^*[\pi^r]$.
\end{theorem}

\begin{proof}
By Lemma~\ref{lem:plan} the plan $\pi^r$ is applicable in $I^*$ and $(I[\pi])^* \subseteq I^*[\pi^r]$. Applying $(\cdot)^*$ to $G \subseteq I[\pi]$ gives $G^* \subseteq (I[\pi])^*$, so
\[
  G^* \subseteq (I[\pi])^* \subseteq I^*[\pi^r]. \qedhere
\]
\end{proof}

\end{document}
